\documentclass{article}
\usepackage{amsmath}
\begin{document}

\section*{Bit Manipulation Techniques}

Bit manipulation allows for efficient data processing and is often used in low-level programming. Here are some common techniques:

\subsection*{Bit Masks}
A bit mask is a binary number used to isolate or modify specific bits in another binary number. Common operations include:

\begin{itemize}
    \item \textbf{Setting a Bit}: To set a specific bit to 1, use the bitwise OR (\texttt{|}) operator with a mask that has the target bit set to 1.
    \begin{verbatim}
    x |= (1 << n); // Set the nth bit of x
    \end{verbatim}
    
    \item \textbf{Clearing a Bit}: To clear a specific bit to 0, use the bitwise AND (\texttt{\&}) operator with the complement of a mask that has the target bit set to 1.
    \begin{verbatim}
    x &= ~(1 << n); // Clear the nth bit of x
    \end{verbatim}
    This can be broken down into the following steps:
    \begin{enumerate}
        \item Shift 1 left by \texttt{n} positions: \texttt{1 << n} shifts the binary 1 to the left by \texttt{n} positions, filling the right with \texttt{n} zeros.
        \item Invert the mask: \texttt{~(1 << n)} inverts all the bits of the result from the previous step.
        \item Apply the bitwise AND operation: \texttt{x \&= ~(1 << n)} applies the bitwise AND operation between \texttt{x} and the inverted mask, clearing the \texttt{n}-th bit of \texttt{x} while leaving all other bits unchanged.
    \end{enumerate}
    
    \item \textbf{Toggling a Bit}: To toggle a specific bit, use the bitwise XOR (\texttt{\^}) operator with a mask that has the target bit set to 1.
    \begin{verbatim}
    x ^= (1 << n); // Toggle the nth bit of x
    \end{verbatim}
    
    \item \textbf{Checking a Bit}: To check if a specific bit is set, use the bitwise AND (\texttt{\&}) operator with a mask that has the target bit set to 1.
    \begin{verbatim}
    bool isSet = x & (1 << n); // Check if the nth bit of x is set
    \end{verbatim}
\end{itemize}

\subsection*{Shifting Bits}
Shifting bits is another common operation used to manipulate data:

\begin{itemize}
    \item \textbf{Left Shift (\texttt{<<})}: Shifts bits to the left, effectively multiplying the number by 2 for each shift position.
    \begin{verbatim}
    x << n; // Shift x left by n positions
    \end{verbatim}
    
    \item \textbf{Right Shift (\texttt{>>})}: Shifts bits to the right, effectively dividing the number by 2 for each shift position.
    \begin{verbatim}
    x >> n; // Shift x right by n positions
    \end{verbatim}
\end{itemize}

\subsection*{One's Complement}
The one's complement operation inverts all the bits in a binary number, turning 1s into 0s and vice versa. This is often used in conjunction with bit masks to clear specific bits.
\begin{verbatim}
~x; // Invert all bits of x
\end{verbatim}

\subsection*{Practical Example}
Let's say you have an 8-bit binary number \texttt{x = 0b10101010} and you want to set the 3rd bit (from the right) to 1:
\begin{verbatim}
x |= (1 << 2); // x becomes 0b10101110
\end{verbatim}

Bit manipulation is a powerful tool that allows for precise control over data at the bit level, which can be very efficient for certain types of operations.

\end{document}
